\section{Síntesis Ejecutiva}

Se propone demostrar que es posible dotar a los asistentes virtuales de un nivel de hiperpersonalización que trascienda la simple recuperación de datos históricos y habilite interacciones contextualizadas, fluidas y centradas en la persona. El problema se delimita en la imposibilidad de los sistemas actuales de recordar la información pertinente de cada usuario de forma estructurada y emplearla en tiempo real para enriquecer la conversación. Sobre esta base, se plantea la hipótesis de que, a partir del análisis de interacciones previas y la ingestión de documentos relevantes, se puede construir una base de datos híbrida —con componentes semánticos y estructurados— capaz de suministrar al asistente el contexto necesario para responder con precisión y pertinencia.

La investigación adopta una metodología experimental dividida en cuatro etapas sucesivas:

\begin{enumerate}
\item Experimentos de \textit{prompting} focalizados
\item Ingestión de documentos para ampliar el contexto
\item Generación de conjuntos de datos sintéticos
\item Alimentación de dichos conjuntos en bases vectoriales y de grafos
\end{enumerate}

Cada fase refina la arquitectura de recuperación y generación de respuestas, combinando recuperación de ambos tipos de bases.

En el desarrollo se implementa un prototipo que integra Redis para el almacenamiento transitorio de sesiones y FAISS para la búsqueda de similitud en espacios vectoriales. A través de sucesivos experimentos —incluidos la desagregación de problemas macro en casos de uso unitarios, la vectorización de mensajes de usuario y la incorporación de marcas temporales— se evidencia que la distancia coseno aplicada tras un filtrado por identificador logra recuperar de forma fiable la información pertinente, aun cuando el corpus crece en complejidad. Esta capacidad es validada, por ejemplo, al responder correctamente preguntas sobre la edad de un familiar o las preferencias de viaje de las hijas del usuario, demostrando que el sistema puede razonar sobre datos cronológicos y semánticos combinados.

Los resultados confirman la viabilidad del enfoque: la combinación de modelos de lenguaje preentrenados y \textit{embeddings} resulta suficiente para procesar y recuperar la información necesaria, siempre que el motor de búsqueda vectorial pueda manejar eficientemente el volumen de datos almacenado. Para su implementación se requiere un servicio de extracción y estructuración de información, una base con capacidad de similitud vectorial y búsqueda híbrida, un servicio de ingestión de datos y otro de generación de respuestas basadas en contexto histórico. Estos componentes pueden desplegarse tanto en infraestructuras locales como en servicios en la nube.

Entre las limitaciones detectadas se señala la necesidad de reconfigurar las fuentes de información cuando se emplean modelos menos potentes o versiones de modelos ya deprecadas, así como la ausencia de un manejo de memoria conversacional que conecte los mensajes de forma continua. Además, el sistema actual no resuelve preguntas sobre múltiples temáticas en una única iteración, lo que impone desafíos de escalabilidad funcional. El trabajo advierte que estos obstáculos deberán afrontarse mediante la integración de servicios de observación paralelos que extraigan y estructuren la información mientras la conversación ocurre.

Como proyección, la tesis plantea que el próximo paso para la adopción masiva de asistentes hiperpersonalizados reside en garantizar la privacidad y el control de los datos mediante tecnologías distribuidas como blockchain, capaces de ofrecer inmutabilidad y gobernanza descentralizada de la información del usuario. Con ello, se aspira a habilitar un ecosistema donde cada individuo comparta —de forma segura y auditable— los fragmentos de su vida digital que considere convenientes, permitiendo a los asistentes virtuales convertirse en verdaderos acompañantes inteligentes, capaces de anticipar necesidades y ofrecer soluciones a medida.

En suma, el trabajo demuestra experimentalmente que la hiperpersonalización en asistentes virtuales no solo es técnicamente factible sino también operativamente viable con recursos accesibles. Sus aportes metodológicos sientan las bases para futuras implementaciones industriales, encaminadas a transformar la experiencia de usuario y a potenciar la eficiencia de los sistemas conversacionales en múltiples dominios.
